% Source-aligned English for printed p.216.
\clearpage
\SourcePage{216}

\LayerTarget{tohoku:proof:coinduced-hom-acyclicity}
\LayerTarget{tohoku:definition:coinduced-module-b-prime}
\noindent make $G$ and $O$ act on it by

\LayerTarget{tohoku:equation:coinduced-u-actions}
\[
(g'f)(g)=f(gg'),
\qquad
(\lambda f)(g)=\lambda^g f(g),
\]
\LayerTarget{tohoku:construction:coinduced-u-module-b-prime}
 one immediately checks that $B'$ thereby becomes a $U$-module and
that an injective $U$-homomorphism $\varphi:B\to B'$ is defined by

\LayerTarget{tohoku:equation:coinduced-u-embedding}
\[
\varphi(b)(g)=g\mathbin{\cdot}b
\]
\LayerTarget{tohoku:reduction:hom-direct-summand-acyclicity}
Thus $B$ embeds in the $U$-module $B'$, and since $B$ is injective, it is a
direct summand of $B'$. Hence $\operatorname{Hom}_O(A,B)$ is (as a
$G$-module) a direct summand of $\operatorname{Hom}_O(A,B')$, and it suffices
to show that the latter is $\Gamma^G$-acyclic.
\LayerTarget{tohoku:definition:conjugation-action-on-hom}
Put $H=\operatorname{Hom}_O(A,B)$ and regard $H$ as a $G$-module by making $G$
act by
$u^g(a)=g\,u(g^{-1}a)$.
One then checks immediately that there is a canonical isomorphism

\LayerTarget{tohoku:equation:hom-coinduction-isomorphism}
\[
\operatorname{Hom}_O(A,B')
=\operatorname{Hom}_O
  \bigl(A,\operatorname{Hom}_{\mathrm Z}(\mathbf Z(G),B)\bigr)
=\operatorname{Hom}_{\mathrm Z}
  \bigl(\mathbf Z(G),\operatorname{Hom}_O(A,B)\bigr)
\]
\LayerTarget{tohoku:construction:two-coinduced-g-actions}
 where the last member
$H''=\operatorname{Hom}_{\mathrm Z}(\mathbf Z(G),H)$ is endowed with the
$G$-module structure defined by
$f^{g'}(g)=g'f(gg')$.
For this, write, for every
$u\in\operatorname{Hom}_O
\bigl(A,\operatorname{Hom}_{\mathrm Z}(\mathbf Z(G),B)\bigr)$,
$u(a)(g)=g\mathbin{\cdot}u_g(a)$. Then the $u_g$ are $O$-homomorphisms from
$A$ to $B$, and $u$ is identified with the family of the $u_g$, itself
identified with an element of
$\operatorname{Hom}_{\mathrm Z}(\mathbf Z(G),H)$. Also consider the
$G$-module $H'$ obtained by putting on
$\operatorname{Hom}_{\mathrm Z}(\mathbf Z(G),H)$ the $G$-module structure
defined by
$(g'f)(g)=f(gg')$.
\LayerTarget{tohoku:equation:coinduced-action-intertwiner}
One obtains a $G$-isomorphism $\psi$ from $H'$ onto $H''$ by putting
$\psi(f)(g)=g^{-1}f(g)$.
\LayerTarget{tohoku:conclusion:coinduced-hom-acyclicity}
It is well known that $H'$ is $\Gamma^G$-acyclic, whatever the $G$-module $H$
may be; hence the same is true of $H''$, q.e.d.

\LayerTarget{tohoku:spectral-sequence:three-equivariant-ext-sequences}
Lemma 5.6.1 allows us to derive from formulas (5.6.5) three spectral sequences
converging to $\operatorname{Ext}_{{\mathbf O},G}(X;A,B)$. To make the initial terms
of the last two explicit, it remains to make explicit the derived functors of
$h_{{\mathbf O},A}$ and $\mathbf h_{{\mathbf O},A}$ regarded as functors on ${\mathbf C}^{{\mathbf O}(G)}$.
\LayerTarget{tohoku:relation:derived-hom-functors}
When they are regarded as functors on the category ${\mathbf C}^{\mathbf O}$ of ${\mathbf O}$-modules
(without operators), their derived functors are, by definition (4.2), the
functors $\operatorname{Ext}_{\mathbf O}^n(X;A,B)$ and, respectively, the functors
$\mathbf{Ext}_{\mathbf O}^n(A,B)$. The same remains true when they are regarded as
functors on ${\mathbf C}^{{\mathbf O}(G)}$, as follows immediately from the definitions and from

\LayerTarget{tohoku:lemma:5.6.2}
\textsc{Lemma 5.6.2.}\quad \textit{If $B$ is an injective
$G$-${\mathbf O}$-module, then it is an injective ${\mathbf O}$-module.}

\LayerTarget{tohoku:proof:lemma-5.6.2}
As in Lemma 5.6.1, we are reduced to the case in which $B$ is the product
sheaf defined by a family $(B_x)_{x\in X}$ of injective $U_x$-modules $B_x$.
It therefore suffices to prove that the $B_x$ are also injective
${\mathbf O}(x)$-modules, which reduces us to the case in which $X$ consists of a single
point $x$.
\LayerTarget{tohoku:reduction:equivariant-injective-underlying}
In fact, the proof does not seem simpler in this case, and the proof we shall
give would apply whenever one has an abelian category ${\mathbf C}$ (such as ${\mathbf C}^{\mathbf O}$) on
which a group $G$ ``acts'', as in\textsuperscript{9)}, permitting passage to
the corresponding category ${\mathbf C}(G)$ (here ${\mathbf C}^{{\mathbf O}(G)}$) of ``objects with an
operator''; assume that ${\mathbf C}$ satisfies $AB$ 5) and has a generator: \textit{then
every injective object $A$ of ${\mathbf C}(G)$ is also injective in ${\mathbf C}$}.
\LayerTarget{tohoku:construction:equivariant-injective-envelope}
To prove this, it plainly suffices to prove that every object $A$ of ${\mathbf C}(G)$
embeds in an object $M$ of ${\mathbf C}(G)$ which is injective in ${\mathbf C}$. For this, it
suffices to examine the construction given in 1.10 for embedding $A$
(regarded as an object of ${\mathbf C}$) in an injective object $M$ of ${\mathbf C}$: one embeds
$A$

\TohokuContinuation{217}
